% Chapter 11: The Erd\H{o}s Distinct Distances Problem
\let\LocalMacro\relax

\chapter{Erd\H{o}s Distinct Distances Problem}

\section{The Problem}

\subsection{Informal}
Scatter a bunch of dots on a piece of paper and look at all the distances between every pair of dots. How many different distances are you guaranteed to see, no matter how you arrange the dots? In 1946, Paul Erd\H{o}s showed that a neat square grid arrangement produces relatively few distinct distances and conjectured that no clever arrangement could do much better. The question stayed open for 64 years until Larry Guth and Nets Katz found a nearly optimal answer using an unexpectedly elegant algebraic trick.

\subsection{Formal}
Erd\H{o}s proved in 1946 that $n$ points in the plane determine at least $cn/\sqrt{\log n}$ distinct distances for some constant $c > 0$, and he conjectured this was essentially tight. Guth and Katz proved in 2010 that the answer is at least $cn/\log n$---nearly matching the upper bound.

\subsection{Historical Context}
Erd\H{o}s was one of the most prolific mathematicians in history, publishing over 1,500 papers. He offered prizes for open problems and could pose beautiful questions on any topic. The distinct distances problem was one of his favorites.

\subsection{What Was Known}
Erd\H{o}s proved in 1946 that a square grid of $n$ points produces about $n/\sqrt{\log n}$ distinct distances, and he conjectured this was essentially optimal. For decades, the best known upper bound was roughly $n/\log^{1/4} n$, which was far from the conjectured answer. Guth and Katz's 2010 breakthrough used the polynomial method---a tool from algebraic geometry---to close the gap almost completely.

\subsection{Why It Matters}
The distinct distances problem reveals fundamental truths about how points can be arranged in space. It connects to sensor networks (how to place sensors to avoid redundant measurements), coding theory (how to pack information efficiently), and the study of geometric patterns in nature.

\subsection{Simplified Version}
If all $n$ points lie on a straight line, there are only $n-1$ distinct distances (the gaps between consecutive points). The question is how much the two-dimensional plane forces more. Erdős showed that $n$ points in the plane must produce at least $c \cdot n/\sqrt{\log n}$ distinct distances, and Guth and Katz proved in 2010 that the answer is at least $c \cdot n/\log n$---nearly matching the upper bound from the square grid.

\section{Mathematical Theory}


\subsection{Prerequisites}
This chapter assumes familiarity with:
\begin{itemize}
\item Basic combinatorics and counting arguments
\item Elementary geometry (distances in the plane)
\item Polynomial algebra (degree, roots)
\end{itemize}

\subsection{The Grid Construction}

\begin{definition}[Integer Grid]
Place $n = m^2$ points at integer coordinates $(i,j)$ for $1 \leq i,j \leq m$. The squared distances are $a^2 + b^2$ for $0 \leq a,b \leq m-1$. The number of distinct squared distances is at most $2m^2 = 2n$, but many of these are not achievable as sums of two squares. A careful analysis shows the number of distinct distances is approximately $n/\sqrt{\log n}$.
\end{definition}

This construction shows $f(n) \leq O(n/\sqrt{\log n})$. The lower bound was the hard part.

\subsection{Previous Bounds}

\begin{itemize}
\item Erd\H{o}s (1946): $f(n) \geq cn/\sqrt{\log n}$ (using the grid construction and counting) \cite{erdos1946}
\item Chung (1984): $f(n) \geq cn^{2/5}$ \cite{chung1984}
\item Clarkson, Edelsbrunner, Guibas, Stolfi, and Welzl (1990): $f(n) \geq cn^{2/3}/\log n$ \cite{clarkson1990}
\item Solymosi and T\'oth (2003): $f(n) \geq cn^{2/3}/(\log n)^{c}$ \cite{solymosi2003}
\end{itemize}

Progress was incremental and frustratingly slow.

\subsection{The Polynomial Method}

The key new idea came from Dvir's proof of the Kakeya conjecture over finite fields (2008), which used polynomials to count geometric configurations \cite{dvir2009}. The insight: if you have many points with few distances, you can construct a nonzero polynomial that vanishes at all the points, and then use algebraic properties to derive a contradiction.

\begin{definition}[Polynomial Method]
Given a set of points $P \subset \mathbb{R}^2$, find a nonzero polynomial $p(x,y)$ of minimum degree that vanishes at all points in $P$. The degree of $p$ constrains the geometry of $P$.
\end{definition}

\section{The Solution}

\subsection{The Big Idea}
Imagine scattering dots on a page and counting all the different distances between pairs. If you arrange them in a neat grid, you get surprisingly few distinct distances---about $n/\sqrt{\log n}$. Erd\H{o}s conjectured this was essentially optimal and spent decades trying to prove that *no clever arrangement could do better*. Guth and Katz's breakthrough was to realize that a tool from *algebraic geometry*---the polynomial method---could attack a problem that seemed purely combinatorial. The core insight: if there are too few distinct distances, you can build a polynomial surface that slices through the points in a way that algebraically *forces* a contradiction. It is like catching a thief by building a cage whose very shape makes escape impossible.

\subsection{What Was Stopping Progress}
For 64 years, all approaches relied on *combinatorial incidence counting*: double-counting arguments, the Szemer\'edi--Trotter theorem, graph-theoretic bounds. These methods could not escape a fundamental bottleneck---they always produced a lower bound with exponent strictly less than 1 because they could not simultaneously control points, lines, and distances without losing a polynomial factor. The $n^{2/3}$ barrier seemed impenetrable. Every technique hit the same wall: pure combinatorics was simply not powerful enough.

\subsection{The Breakthrough}
The key idea came from a completely different field. In 2008, Zeev Dvir proved the finite-field Kakeya conjecture using *polynomials*: construct a nonzero polynomial of controlled degree that vanishes on the set, then use the Schwartz--Zippel lemma to derive a contradiction. Guth and Katz realized this algebraic technique could be adapted to the real plane via *polynomial partitioning*: divide $\mathbb{R}^2$ into cells using a polynomial's zero set, control the points in each cell, and use a *joints argument* to count intersections. This transforms a geometric problem into an algebraic one---and escapes the $n^{2/3}$ barrier.

\subsection{How It Works}

Erd\H{o}s proved in 1946 that $n$ points in the plane determine at least $cn/\sqrt{\log n}$ distinct distances, and his square grid construction showed that this was essentially tight---so the answer should be $\Theta(n/\sqrt{\log n})$. But after that promising start, progress stalled for decades. The best lower bound achieved by 2003 (Solymosi--T\'oth) was only $cn^{2/3}/(\log n)^c$ \cite{solymosi2003}. The exponent $2/3$ was nowhere near the $1$ needed to match the upper bound. The gap between what was known ($n^{2/3}$) and what was conjectured ($n/\sqrt{\log n}$) was enormous, and it had persisted in some form since the 1980s.

All prior approaches relied on \emph{combinatorial incidence counting}: lines, circles, and distances were studied through double-counting arguments, the Szemer\'edi--Trotter theorem, and graph-theoretic bounds. These methods could not escape a fundamental bottleneck. When you try to count how many points lie on how many lines, the best tools give you bounds of the form $n^{2/3}$ (from incidence geometry) or $n/\sqrt{\log n}$ (from the grid construction). The combinatorial methods always produced a lower bound with exponent strictly less than 1 because they could not simultaneously control all three variables---points, lines, and distances---without losing a polynomial factor. Every known technique hit the same wall: the $n^{2/3}$ barrier seemed impenetrable.

The breakthrough came from a completely different field. In 2008, Zeev Dvir proved the finite-field Kakeya conjecture using \emph{polynomials}: he showed that if a set contains a line in every direction over a finite field, it must be large, by constructing a nonzero polynomial of controlled degree that vanishes on the set and then using the Schwartz--Zippel lemma \cite{dvir2009}.

Guth and Katz realized that this algebraic technique could be adapted to the real plane. The key innovation was \textbf{polynomial partitioning}: given $n$ points, find a nonzero polynomial $p(x,y)$ of degree $d$ such that the zero set of $p$ divides $\mathbb{R}^2$ into $O(d^2)$ cells, each containing at most $n/d^2$ points. This transforms a geometric problem into an algebraic one. Inside each cell, the points are well-separated, and across cells, the polynomial structure constrains how distances can repeat.

The second key idea was a \emph{joints argument}: if many lines meet in a point, that point is a ``joint'' where three non-coplanar lines intersect. Guth and Katz showed that a polynomial of degree $d$ can have at most $O(d^2)$ joints per line, giving a new incidence bound that escaped the $n^{2/3}$ barrier.

Here is how to think about the key insight. Imagine you have $n$ points scattered in the plane and very few distinct distances. That means many pairs of points are the same distance apart, so many pairs ``agree'' on a distance. Guth and Katz's idea is to build a polynomial surface (like a wavy sheet) that slices through the plane and puts roughly equal numbers of points in each piece. Because the polynomial has algebraic structure, it controls how the points can be arranged: you cannot pack too many points into a small region without the polynomial's degree forcing a contradiction. The polynomial acts like a ``net'' that catches any configuration with too few distances.

More concretely: if there were far fewer than $n/\log n$ distinct distances, then some distance $d$ would appear many times. That means many points lie on many circles of radius $d$. The polynomial partitioning forces these circles to intersect the polynomial's zero set in controlled ways, and the joints argument counts those intersections. The algebraic structure of the polynomial---its degree, its factorization, its behavior at special points---gives you counting power that pure combinatorics never could.

The result is $f(n) \ge cn/\log n$, which almost matches the grid construction's upper bound of $O(n/\sqrt{\log n})$. The $\log n$ gap remains, but the polynomial method closed a gap that had been open for 64 years.

Larry Guth and Nets Katz proved in 2010 \cite{guth2015}:

\begin{theorem}[Guth--Katz, 2010 \cite{guth2015}]
For any set of $n$ points in the plane, the number of distinct distances is at least $cn/\log n$ for some constant $c > 0$.
\end{theorem}

This was essentially optimal (matching Erd\H{o}s's grid construction up to the $\log n$ factor).

Their proof combined several ideas:

\begin{enumerate}
\item \textbf{Polynomial partitioning}: Divide the plane into regions using a polynomial surface, then study how points and distances distribute across regions.
\item \textbf{Incidence geometry}: Count how many times points lie on lines (or curves) using the Szemer\'edi--Trotter theorem.
\item \textbf{Joints argument}: Count the number of ``joints'' (points where three non-coplanar lines meet) to control the geometry.
\end{enumerate}

\begin{highlightbox}
The Guth--Katz proof was a masterpiece of combining techniques from different areas. The polynomial method, originally from finite fields, was adapted to the real plane using a novel partitioning argument.
\end{highlightbox}

\subsection{Subsequent Developments}

\begin{itemize}
\item \textbf{Improved bounds}: Later work by Dong, Guth, and others improved the bound slightly.
\item \textbf{Higher dimensions}: The polynomial method has been extended to distinct distances in $\mathbb{R}^d$ for $d \geq 3$.
\item \textbf{Other metric spaces}: The polynomial method has been applied to distinct distances in various settings.
\end{itemize}

\section{Impact}

\begin{itemize}
\item \textbf{Polynomial method}: This proof cemented the polynomial method as a major tool in combinatorial geometry.
\item \textbf{Incidence geometry}: The Guth--Katz approach has influenced many problems in incidence geometry.
\item \textbf{Interdisciplinary connections}: The techniques connect algebra, geometry, and combinatorics in new ways.
\end{itemize}

\section{Exercises}

\begin{exercise}[Basic]
Show that $n$ collinear points determine at most $n-1$ distinct distances. Why doesn't this contradict the theorem?
\end{exercise}

\begin{exercise}[Intermediate]
Construct a set of $n$ points on a circle and count the number of distinct distances. How does this compare to the grid construction?
\end{exercise}

\begin{exercise}[Challenge]
Explain why the polynomial method works better in finite fields than in $\mathbb{R}$. What new difficulties arise in the real case?
\end{exercise}


\begin{exercise}[Computational]
Write a Python/SageMath script to explore a concept from this chapter. See the appendix for starter code.
\end{exercise}

\section{Timeline}

\begin{tabularx}{\linewidth}{lX}
\toprule
\textbf{Year} & \textbf{Event} \\ \midrule
1946 & Erd\H{o}s poses the distinct distances problem \cite{erdos1946} \\ 
1984 & Chung improves the lower bound to $n^{2/5}$ \cite{chung1984} \\ 
2008 & Dvir proves the finite field Kakeya conjecture using polynomials \cite{dvir2009} \\ 
2010 & Guth and Katz prove $f(n) \geq cn/\log n$ \cite{guth2015} \\ 
2015+ & Further improvements and extensions \\ \bottomrule
\end{tabularx}

\section{Citations}

\begin{enumerate}
\item Guth, L., \& Katz, N. (2015). ``On the Erd\H{o}s distinct distances problem in the plane.'' Annals of Mathematics, 181(1), 155--190.
\item Erd\H{o}s, P. (1946). ``On sets of distances of $n$ points.'' American Mathematical Monthly, 53(5), 248--250.
\item Dvir, Z. (2009). ``On the size of Kakeya sets in finite fields.'' Journal of the AMS, 23(4), 1107--1119.
\item Chung, F. (1984). ``The number of different distances determined by $n$ points in the plane.'' Journal of Combinatorial Theory, Series A, 36(3), 342--354.
\item Clarkson, K. L., Edelsbrunner, H., Guibas, L. J., Stolfi, M., \& Welzl, E. (1990). ``Combinatorial complexity bounds for arrangements of curves and spheres.'' Discrete \& Computational Geometry, 5(2), 99--160.
\item Solymosi, J., \& T\'oth, C. D. (2003). ``Distinct distances in the plane.'' Discrete \& Computational Geometry, 30(3), 433--448.
\end{enumerate}
